This long exact sequence is simply the restriction to
$\mathbb V(\underline E)^*$ of the Koszul resolution of
$\underline O_{e(S)}$ on $\mathbb V(\underline E)$.

Push the long exact sequence (1.1.5) forward to $\mathbb P(\underline E)$ by
the affine morphism $q$.  By (1.1.1) and the formula
$q_*{a'}^*=q_*q^*p^*$, one obtains a long exact sequence

\begin{equation}\tag{1.1.6}
\cdots\longrightarrow
\bigoplus_{n\in\mathbb Z}p^*\bigwedge^i\underline E(n)
\xrightarrow{\ \iota_X\ }
\bigoplus_{n\in\mathbb Z}p^*\bigwedge^{i-1}\underline E(n)
\longrightarrow\cdots .
\end{equation}

It splits into short exact sequences, the images of (1.1.3),

\begin{equation}\tag{1.1.7}
0\longrightarrow
\bigoplus_{n\in\mathbb Z}\Omega^i_{\mathbb P(\underline E)/S}(n)
\longrightarrow
\bigoplus_{n\in\mathbb Z}p^*\bigwedge^i\underline E(n)
\xrightarrow{\ \iota_X\ }
\bigoplus_{n\in\mathbb Z}\Omega^{i-1}_{\mathbb P(\underline E)/S}(n)
\longrightarrow0.
\end{equation}

The exact sequences (1.1.7) are direct sums of the exact sequences

\begin{equation}\tag{1.1.8}
0\longrightarrow\Omega^i_{\mathbb P(\underline E)/S}(n)
\longrightarrow p^*\bigwedge^i\underline E(n-i)
\xrightarrow{\ \iota_X\ }
\Omega^{i-1}_{\mathbb P(\underline E)/S}(n)
\longrightarrow0.
\end{equation}

Recall (EGA III, 2.1.15, or FAC) that the sheaves
$R^ip_*\underline O(k)$ are locally free and their formation is compatible
with arbitrary base change, and that

\begin{equation}\tag{1.1.9}
R^ip_*\underline O(k)=0
\quad\text{if }i\ne0,r,
\quad\text{or if }i=0, k<0,
\quad\text{or if }i=r, k>-r-1,
\end{equation}

and hence that $R^ip_*p^*\bigwedge^j\underline E(k)$ vanishes under the
same conditions.

The exact sequences (1.1.8) give a resolution

\begin{equation}\tag{1.1.10}
0\longrightarrow\Omega^i_{\mathbb P(\underline E)/S}(n)
\xrightarrow{\ E\ }p^*\bigwedge^i\underline E(n-i)
\xrightarrow{\ \iota_X\ }\cdots\longrightarrow
p^*\bigwedge^0\underline E(n)\longrightarrow0
\end{equation}

of $\Omega^i_{\mathbb P(\underline E)/S}(n)$.  The spectral sequence
defined by this resolution degenerates ($E_1=E_\infty$) for degree reasons
by (1.1.9); for $0<j<r$, it gives

\begin{equation}\tag{1.1.11}
R^jp_*\Omega^i_{\mathbb P(\underline E)/S}(n)
=\mathcal H\left(\cdots\xrightarrow{\ \iota_X\ }
p_*\bigwedge^{i-j}\underline E(n-i+j)
\xrightarrow{\ \iota_X\ }\cdots\right).
\end{equation}

Since $r\geq1$, a depth argument shows that

\begin{equation}\tag{1.1.12}
\bigoplus_{n\in\mathbb Z}p_*p^*\bigwedge^i\underline E(n)
=p_*q_*q^*p^*\bigwedge^i\underline E
={a'}_*{a'}^*\bigwedge^i\underline E
=a_*a^*\bigwedge^i\underline E.
\end{equation}

%%%%%%%%%% source scan idx 50 | folio 43 %%%%%%%%%%

Thus $p_*p^*\bigwedge^i\underline E(n)$ is the homogeneous part of weight
$n+i$ in $a_*a^*\bigwedge^i\underline E$, namely
$\operatorname{Sym}^n(\underline E)\otimes\bigwedge^i\underline E$.
Remembering that (1.1.5) comes from the Koszul resolution, (1.1.11) shows
that, for $0<j<r$, the sheaf
$R^jp_*\Omega^i_{\mathbb P(\underline E)/S}(n)$ can be nonzero only when
$i-j=0$ and $n-i+j=0$, that is, only when $n=0$ and $i=j$.  In that case,
moreover, $R^ip_*\Omega^i_{\mathbb P(\underline E)/S}$ is isomorphic to
$\underline O_S$.

For $j=0$, (1.1.10) also gives

\begin{equation}\tag{1.1.13}
p_*\Omega^i_{\mathbb P(\underline E)/S}(n)
=\ker\!\left(\iota_X:
p_*\bigwedge^i\underline E(n-i)\longrightarrow
p_*\bigwedge^{i-1}\underline E(n-i+1)\right),
\end{equation}

and the same argument shows that $\iota_X$ is injective when $n-i\leq0$.

By (1.1.8), with $n=0$ and $i=r+1$, one has

\begin{equation*}
\Omega^r_{\mathbb P(\underline E)/S}\simeq
p^*\bigwedge^{r+1}\underline E(-r-1).
\end{equation*}

When $S$ is the spectrum of a field, the vanishing assertions in 1.1(ii)
and (iii) now follow from the preceding argument for $j\ne r$, and follow
from it by duality ([2], V, 11.2, p.~211) for $j=r$.  In particular, for
any $j$ and $n$ there is at most one $i$ for which
$R^ip_*\Omega^j_{\mathbb P(\underline E)/S}(n)\ne0$.  This proves (i),
since $\Omega^j_{\mathbb P(\underline E)/S}$ is flat over $S$ (the hardest
part is locating the reference in EGA III; we suggest 6.10.5, 7.4.1, and
7.5.5, followed by a passage to the limit).

Knowing that $R^ip_*\Omega^i_{\mathbb P(\underline E)/S}$ is an invertible
sheaf for $0\leq i\leq r$, it remains only to prove that, when $S$ is the
spectrum of a field, $\eta^i\ne0$.  Indeed, $\eta^r$ is the cohomology
class of a rational point (the intersection of $r$ hyperplanes), hence a
basis of $H^r(\mathbb P^r,\Omega^r)$ (FGA, Exposé 149, no.~4).

1.2. The Euler--Poincaré characteristics of the $\Omega^i(n)$ may be
calculated directly from (1.1.10) or (1.1.13).  The rather unappetizing
result is

\begin{equation*}
\chi\bigl(\mathbb P^r_k,\Omega^j_{\mathbb P^r_k}(n)\bigr)
=\sum_{0\leq i\leq j}(-1)^{i-j}
\binom{r+1}{i}\binom{r+n-i}{r}.
\end{equation*}

